\section{Conclusion}
\label{sec:conclusion}

This study presents the \textit{MVC Test Orchestrator}, a framework
designed to detect selected MVC conformance violations in newly generated Python artifacts
within LLM-assisted software development pipelines. By distributing functional responsibilities across specialized LLM-based agents and rule-based software modules, each operating within clearly bounded task domains, the current system aims to enable structured control over complex architectural design tasks.

The empirical effectiveness of the system was validated through 
experiments conducted across a dataset of 20 distinct Software 
Requirements Specifications (SRS). The experimental outcomes demonstrate
that the system achieved a macro-average BERTScore of 0.808 and an average extraction score (ACES) of \textbf{74.0\%}
across Model, View, and Controller components, alongside an architectural
violation-detection recall (VDR) of \textbf{98.8\%} (96.2\% precision, ACP) and an
autonomous remediation success rate (RSDA) of \textbf{92.6\%} among detected
violations (91.6\% end-to-end, RSR\textsubscript{e2e}). The five-expert evaluation panel exhibited excellent
inter-rater agreement (ICC(A,5)\,$=$\,0.942), supporting the reliability of the reported ACES scores, and
the multi-iteration analysis yielded a macro-average Consistency Score (CS) of 70.6\% across three
independent extraction runs. As expected for a deterministic,
rule-based transformation, the $0\%$ Scaffolder
Error Rate (SER) documented during the \texttt{/scaffold} phase confirms
that the rule-based scaffolding stage faithfully translated the
Architectural Map into the file-system structure across all 20 SRS
documents, consistent with its intended role as a rule-based, non-LLM
safety perimeter (Section~\ref{sec:phase3}).

A key contribution of the \textit{MVC Test Orchestrator} lies in its synthesis of
Retrieval-Augmented Generation (RAG) paradigms with rule-based Abstract
Syntax Tree (AST) analysis. This integration establishes an end-to-end
traceability chain spanning from natural language
requirements to the executable codebase. The pipeline preserves developer oversight through a Human-in-the-Loop (HITL) integration at key decision points. Collectively, these results support the feasibility of the tested pipeline for selected MVC conformance tasks within newly generated Python artifacts, but do not establish superiority, longitudinal prevention of architecture erosion, functional correctness, or generalization beyond the studied model, corpus, and Python rule set. The 
limitations of this study, including dataset scope and model-specific 
constraints, are discussed in Section~\ref{sec:limitations}; directions 
for future work are outlined in Section~\ref{sec:future}.